\documentclass{report}
\usepackage{amsmath,amssymb}
\setlength{\parindent}{0mm}
\setlength{\parskip}{1em}
\begin{document}
\begin{center}
\rule{6in}{1pt} \
{\large Christopher A. Leibs \\
{\bf Least-Squares Finite Element Method and Nested Iteration for Electromagnetic Two-Fluid Plasma Models}}

526 UCB \\ University of Colorado \\ Boulder \\ CO 80309
\\
{\tt leibs@colorado.edu}\\
Thomas Manteuffel\end{center}

Efforts are currently being directed towards a fully implicit,
electromagnetic, JFNK-based solver, motivating the necessity of
developing a fluid-based, electromagnetic, preconditioning strategy [1].
The two-fluid plasma (TFP) model is an ideal approximation to the kinetic
Jacobian. The TFP model couples both an ion and an electron fluid with
Maxwell's equations. The fluid equations consist of the conservation of
momentum and number density. A Darwin approximation of Maxwell is used to
eliminate light waves from the model in order to facilitate coupling to
non-relativistic particle models. We analyze the TFP-Darwin system in the
context of a stand-alone solver with consideration of preconditioning a
kinetic-JFNK approach.

The TFP-Darwin system is addressed numerically by use of nested iteration
(NI) and a First-Order Systems Least Squares (FOSLS) discretization. An
important goal of NI is to produce an approximation that is within the
basin of attraction for Newton's method on a relatively coarse mesh and,
thus, on all subsequent meshes. After scaling and modification, the
TFP-Darwin model yields a nonlinear, first-order system of equations
whose Fr\'{e}chet derivative is shown to be uniformly
$\mathcal{H}^1$-elliptic in a neighborhood of the exact solution.
$\mathcal{H}^1$ ellipticity yields optimal finite element performance and
linear systems amenable to solution with Algebraic Multigrid (AMG). To
efficiently focus computational resources, an adaptive mesh refinement
scheme, based on the accuracy per computational cost, is leveraged.
Numerical tests, including magnetic reconnection, demonstrate the
efficacy of the approach, yielding an approximate solution within
discretization error in a relatively small number of computational work
units.
\\\\
\noindent[1] G. Chen, L. Chac\'{o}n, Computer Physics Communications, 185(10), 2014.


\end{document}
